\section{Objetivo Geral}

Desenvolver um sistema de acompanhamento da composição corporal baseado em bioimpedância
elétrica e inteligência artificial, capaz de integrar os dados de uma balança de oito
eletrodos, interpretar os resultados automaticamente e fornecer informações e
recomendações personalizadas ao usuário.

\section{Objetivos Específicos}

\begin{itemize}
    \item Desenvolver um protótipo funcional (Prova de Conceito, PoC) de arquitetura de
    software, englobando \emph{backend} e interface web progressiva (PWA), para o
    gerenciamento seguro dos dados.
    \item Integrar o protótipo aos dados da balança CF577BLE por meio de comunicação
    Bluetooth Low Energy (BLE) e das APIs do fabricante.
    \item Implementar mecanismos básicos de anamnese para registrar fatores contextuais
    (hábitos, hidratação) que influenciam as medições de bioimpedância.
    \item Explorar e desenvolver um modelo computacional inicial para ajuste das medições
    da balança, visando aproximá-las de padrões validados para a população brasileira.
    \item Integrar e especializar um \emph{Small Language Model} (SLM) para interpretação
    dos dados de composição corporal por meio de regras de interpretação, engenharia de
    \emph{prompt} e integração com modelos biométricos fundamentados na literatura.
    \item Avaliar o desempenho técnico do sistema e a coerência das respostas geradas pelo
    SLM.
    \item Realizar testes preliminares de usabilidade com uma amostra reduzida de usuários
    voluntários, avaliando a aceitação e a viabilidade da solução.
    \item Implementar uma interface conversacional via WhatsApp para interação com o
    assistente virtual.
    \item Criar uma solução de \emph{Machine Learning} (ML) para o aprendizado
    personalizado dos dados biométricos de cada usuário ao longo do tempo.
\end{itemize}

Para atingir esses objetivos, foi desenvolvido o sistema \sintonia, composto por uma
balança de bioimpedância integrada via Bluetooth, um servidor de aplicação responsável
pelo processamento dos dados, uma interface web progressiva para interação com o usuário e
uma camada de inteligência artificial voltada à interpretação dos resultados. Os detalhes
da arquitetura, das tecnologias empregadas e dos mecanismos de processamento são
apresentados nos capítulos seguintes.
